\section{Grouped Euler minima and progress-driven screening}
\label{sec:euler-aggregate}

The source Euler compiler of Section~\ref{sec:window-euler} evaluates
canonical Euler characteristic without constructing a normal vector.
Its first implementation still scans all physical corners at every tested
Q corner. This continuation groups equal projected forms within each
source vertex and compiles one small nonlinear objective for the remaining
corner queries. It retains the cheap one-dimensional branch and builds
no grouping when the first corner is already positive.

\subsection{Anchor shifts and constant vertex groups}

For the updated basis $K'$ write
\[
 \alpha_a=\sum_j\gamma_jK'_{a,j},\qquad
 F(z)=\alpha\cdot z-\sum_v\lambda_v\min_{c\in G_v}\ell_c(z).
\]
Choose one actual form $a_v$ as the anchor for vertex $v$, and let $D_v$
be the distinct nonzero differences $\ell_c-a_v$. Define
\[
 \bar\alpha=\alpha-\sum_v\lambda_va_v.
\]
Then
\[
 F(z)=\bar\alpha\cdot z-
       \sum_{v:D_v\ne\varnothing}\lambda_v
                   \min\!\left(0,\min_{f\in D_v}f\cdot z\right).
\]
The zero anchor difference is supplied explicitly. If a vertex group has
only one projected form, it contributes only to $\bar\alpha$ and disappears
from the nonlinear minima. Adding any common linear offset to a vertex's
forms, together with its exact link-weighted change to the signed linear
term, cancels under this normalization.

\begin{proposition}[Grouped Euler objective is exact]
The grouped functional equals canonical source Euler characteristic at
every parameter vector. Duplicated forms and constant vertex groups can
be omitted as above without changing any corner value or exclusion result.
\end{proposition}
\begin{proof}
For each vertex,
$\min_c\ell_c(z)=a_v\cdot z+\min(0,\min_{f\in D_v}f\cdot z)$.
Substitution and collecting the anchor terms give the displayed identity.
Duplicates do not alter a minimum. An empty difference set contributes
only its anchor linear form. The equality is algebraic and does not require
positivity or integral coordinates. Applied to the normalized nonnegative
Q section, it therefore retains the preceding complete nonpositive-corner
criterion and its independent source interpretation.
\end{proof}

\subsection{Only refine the objective after a nonpositive first corner}

The implementation tests the first Q corner with the existing direct
canonical functional. A positive value immediately resumes ordinary
standard discovery. A point section also finishes without aggregation.
Only resuming to another corner triggers compilation of distinct forms
and the anchor-adjusted linear term. Later corners use that one compiled
objective. Empty sections and the scalar nullity-one branch keep their
previous behavior; higher-nullity generic fallbacks remain available.

The grouping belongs to one updated support, basis and source projection.
It is not reused across different sectors. Equality is exact rational
value equality, with no numerical tolerance. Compilation completes before
a grouped value is used; callback interruptions propagate and cannot give
an exclusion from a partial list of corners. Positive corner values still
prove neither connectedness nor disc essentiality.

\subsection{Distinct-form bound and arithmetic cost}

The source support contraction of Section~\ref{sec:incoming-sectors} removes matching
edges with identically zero allowed-Q labels. Corner forms are constant
along every such contracted component. Equal-form grouping can merge those
components further, including dropping groups that become constant; it
cannot split them. Consequently the total number $p$ of distinct forms in
nonconstant vertex groups is bounded by the original retained class bound
$p\le8k$. This bound is about actual source matching geometry, not arbitrary
families of affine functions.

Compiling $\alpha$ costs $O(kd)$ arithmetic operations. Scanning the physical
corner forms, adjusting anchors and creating distinct differences costs
$O(Td+pd)$ coefficient operations, plus exact grouping/index costs.
Each later corner evaluates only $O(pd)$ coefficients. Including the
unaggregated first corner and compilation, $m$ tested corners cost
$O(d(T+k+mp))$ arithmetic operations after source projection. At nullity
at most three, the normalized polygon has $O(k)$ vertices; thus this
implements the proposed $O(T+k^2)$ post-projection coefficient arithmetic.
The normalized chart and polygon costs are still charged.

Python dictionaries implement exact form grouping. Their expected lookup
cost and rational hashing/bit costs are distinct from the scalar-operation
bound; no deterministic linear-time hash guarantee is asserted.
Source preparation, potential projection, coefficient storage, actual
ray enumeration, full coordinate output and independent proof replay also
remain separately charged. This improves a local polynomial factor. A
global centre family, logarithmic distance bound and general quasi-polynomial
unknot recognition theorem remain unproved.

\subsection{Source comparisons and timings}
All 1,417 maintained tests pass in 254.597 seconds; the fourteen focused
tests pass in 22.808 seconds. Every tested small source window compares
the grouped objective, direct canonical formula and fresh source Euler data,
including scaling by $2^{90}+1$. Genuine source tests check anchor-gauge
invariance, constant-group elimination and the $8k$ distinct-form bound.
A positive first corner is tested with aggregation disabled, proving that
no unused objective is compiled. Higher-nullity fallback behavior is retained.

The 85-source audit exactly preserves every disabled output, every enabled
verdict and all 31 native positives. Seven ordinary window discs receive
fresh Regina confirmation. The twenty-two capped cases and thirty-two
completed local misses remain unchanged. The motivating genus-one case
builds seven grouped objectives with 127 retained distinct forms in total;
trefoil and figure-eight build ten and twelve, with totals 259 and 289.
Guard work rises from 106,508 to 108,245, from 171,166 to 173,683 and from
300,945 to 304,918 on these three cases. Grouping preparation is real work,
not a universal guard-count improvement. All 605 captured runtime pins stay
fixed. The 185.301-second audit is not a paired throughput comparison.

An initial three-round run completes 120 measured calls and forty warmups.
The final five-round comparison completes 200 measured calls and forty
warmups, freezing four preceding-release modules at \code{2c2db6b42}:
Euler screening, residual queries, seed queries and recognition. Local imports
are redirected only to corresponding frozen dependencies; no timed patch
context is charged. Both arms finish their configured native work and
agree on statuses and tested positive certificate digests. Native timings
include source and grouping preparation, actual queries, independent successful
replay and serialized output. The pilot and final records are both retained.

\begin{samepage}
\noindent\textbf{Complete native queries with radius-two windows.}
\begin{center}
\small
\begin{tabular}{lrrrrr}
\toprule
Input & previous ms & current ms & old/new & old A/A & new A/A\\
\midrule
empty circle & 5.806 & 5.819 & 0.991 & 0.998 & 1.002\\
optimized positive & 57.007 & 57.541 & 1.001 & 0.989 & 0.982\\
genus one miss & 190.226 & 178.944 & 1.064 & 1.014 & 0.987\\
trefoil & 314.882 & 306.607 & 1.027 & 1.003 & 1.030\\
figure eight & 540.842 & 522.980 & 1.034 & 0.982 & 1.004\\
\bottomrule
\end{tabular}
\end{center}

\end{samepage}

The native genus-one, trefoil and figure-eight first-run gains are about
1.064, 1.027 and 1.034. Early controls that return before aggregation remain
close to parity. These are small-case measurements over the preceding
verifier-context implementation, not comparisons with omitting windows.

\begin{samepage}
\noindent\textbf{Complete recognition with radius-two windows enabled.}
\begin{center}
\small
\begin{tabular}{lrrrrr}
\toprule
Input & previous ms & current ms & old/new & old A/A & new A/A\\
\midrule
empty circle & 0.018 & 0.017 & 1.023 & 0.966 & 0.961\\
optimized positive & 53.087 & 53.316 & 0.996 & 1.012 & 0.997\\
genus one miss & 199.700 & 193.789 & 1.037 & 1.027 & 1.025\\
trefoil & 335.007 & 326.341 & 1.042 & 0.992 & 0.991\\
figure eight & 564.884 & 546.939 & 1.031 & 0.989 & 0.990\\
\bottomrule
\end{tabular}
\end{center}

\end{samepage}

The first complete-recognition gains are 1.037, 1.042 and 1.031 on the
same cases. Because these are modest and the genus-one A/A controls are
about 1.027 and 1.025, a separate fifteen-round comparison repeats the
genus-one and figure-eight cases. It completes 120 measured calls and eight
warmups with complete recognition and serialized output. The old/new ratios
are 1.002 and 1.000, approximately parity; A/A controls are 1.008 and 1.068
on the genus-one case and 0.991 and 1.026 on the figure-eight. The repeat
therefore does not support a reliable small-case wall-clock gain. Both
records are published. No further ratio is inferred from the microsecond
empty-recognition control.

The concrete result is the exact grouped objective and improved
post-projection coefficient bound, not a uniform practical speedup.
Preparation can outweigh saved scans on short or small queries.
Radius zero remains much faster on the small controls; the window stage
continues to default off. A complete local window still supplies no general
negative knot certificate or quasi-polynomial coverage theorem.

\begin{sloppypar}
Runtime release \code{71e1e5225}, tests, the source audit, pilot, final timing
run and longer repeat are retained under \code{data/euler-aggregate-*}.
From \code{fast/}, use \code{normal\_orbit\_research.euler\_aggregate} with
\code{audit --fresh-regina} or \code{benchmark --rounds 5}.
Publication checks verify 605 frozen/current runtime pins, 604 preceding
release evidence pins, all timing outcomes and all 31 saved positives
with grouping, the batch verifier context, prepared observers, cached modes,
Euler compilation, matching updates, residual planning, coherent extraction,
ray enumeration and orbit producers disabled. Eight saved generic legacy
proofs also receive producer-disabled fresh/context replay.
\end{sloppypar}

